\cotefile{57}
\scanpage{65}
Plan
% archive sheet number: 64

\scanpage{66}
\uline{Picard}

1. Théorèmes d'existence absolus

2. Théorèmes d'existence relatifs.

\noindent\llap{$\left\{\vphantom{\rule{0pt}{2em}}\right.$\enspace}3. \inserted{\fbox{Critères d'équivalence et}} Théorèmes de finitude. [\uline{Néron}-\uline{Igusa} exclus ??]

4. $\underline{\mathrm{Pic}}^{\tau}_{X/S}$, équivalence numérique, gpe de Néron Sévéri.

\deleted{\unsure{5|6}}6. Critères de \deleted{\unsure{propreté}} séparation, de propreté pour $\underline{\mathrm{Pic}}_{X/S}$.

\deleted{\illegible{1}} \inserted{\fbox{Critères de simplicité pour $\underline{\mathrm{Pic}}_{X/S}$}}.

\deleted{\unsure{7}}8. Albanese

\deleted{\unsure{8|7}}7. Dualité schémas abéliens.
% slanted; a long curved leader runs up from the left margin to this group.
\margin{\unsure{utiliser et|utiliser les} \unsure{rester sur|variétés des} les schémas en groupes\ldots}

9. \uline{Picard} des courbes relatives [\inserted{\fbox{notamment \unsure{autodualité|auto-dualité}.}} Torelli ?]

\deleted{6}5 \deleted{\unsure{Th.}} Variations, sur \unsure{le|les|la} \unsure{carré|carrés|cour} et le cube [\unsure{multiplicativité} du $\underline{\mathrm{Pic}}^{\tau}_{X/S}$]
\inserted{\uwave{Correspondances divisorielles}}

10. Compléments au dessus d'un corps.

[notamment cas des variétés non complètes, des diviseurs non de Cartier, le point de vue birationnel ; \deleted{\unsure{etc.]}} Critères d'équivalence dans le cas normal etc \deleted{\unsure{etc]}} ; aussi rappel des résultats spéciaux obtenus dans les paragraphes antérieurs]

11. Phénomènes spéciaux en car. 0.

12. (?) Picard locaux et progroupes \unsure{---|\ldots}

\diagram{\begin{tikzcd}[ampersand replacement=\&,row sep=small,column sep=small]
 \& 1 \arrow[dl] \arrow[dr,Rightarrow] \arrow[ddd,dashed] \& 2 \arrow[d,Rightarrow] \\
 {} \& {} \& 3 \arrow[d,Rightarrow] \\
 5 \& {} \& 4 \arrow[d,Rightarrow] \arrow[dl,dashed] \\
 {} \& 6 \arrow[ddr,Rightarrow] \& 5 \arrow[dd,Rightarrow] \\
 {} \& {} \& {} \\
 {} \& {} \& 7 \arrow[d,Rightarrow] \\
 {} \& {} \& 8 \arrow[d,Rightarrow] \\
 {} \& {} \& \unsure{9}\,\deleted{\unsure{10|11}} \arrow[r] \& 10
\end{tikzcd}}
% archive sheet number: 65

\scanpage{67}
\uline{Th. de finitude}
% archive sheet number: 66

\scanpage{68}
\fbox{(\uline{Dém} de Matsusaka)}

% slanted
\margin{finitude de $\underline{\mathrm{Pic}}^{\tau}_{X/S}$}
Soit $X$ une surface projective non singulière sur un corps algébriquement clos $k$. \inserted{Soit $H$ un diviseur ample.} Rappelons la formule de R-R. faible
\[
\chi(D)=\frac12D(D-K)+c
\]
où $c$ est une constante. \deleted{Soit $D'$ numériquement équivalent à $0$, on trouve} \inserted{Comme}
\[
\chi(D)=\ell(D)+\ell(K-D)-\dim H^1(\unsure{x|X},\unsure{L|\mathcal{L}}(D))\qquad\text{d'où}
\]
\[
\begin{aligned}
&\deleted{\chi(}\\
&\deleted{\chi(D)}\\
&\ell(D)+\ell\unsure{(K-D))|(K-D)}\geq\frac12D(D-K)+c
\end{aligned}
\]

\deleted{Donc on en conclut}

Supposons p.ex. que $\ell(K-D)=0$, ce qui est le cas si $H.(K-D)\leq0$ i.e. si
\[
\text{a)}\qquad \deleted{\unsure{HK<DH|H(K-D)}} (D-K).H\geq0
\]
Alors on a $\ell(D)\geq\frac12D(D-K)+c$, donc a) et \deleted{\unsure{si on}}
\[
\text{b)}\qquad \frac12D(D-K)+c\geq0
\]
implique $\ell(D)>0$. Les conditions a) et b) sont satisfaites p.ex. pour \deleted{\illegible{1}} $D=nH$, $n$ assez grand. Si alors $\Delta$ \deleted{\unsure{re}} est un diviseur numériquement équivalent à $0$, les conditions a) et b) sont vérifiées pour $\Delta$, donc on a
\[
\ell(D+\Delta)>0
\]
donc $D+\Delta$ est représenté par un diviseur positif. Donc :
% archive sheet number: 67

\scanpage{69}
\uline{Proposition} Soit $D$ un diviseur satisfaisant a) et b), et soit $M$ la famille des \uline{diviseurs $>0$} \uline{numériquement} équivalents à $D$. Alors tout diviseur \unsure{numériqt|numériquement} équivalent à $0$ est lin. équivalent à un diviseur de la forme $D'-D$, avec $D'\in M$, \deleted{\unsure{donc}} \unsure{i.e|i.e.} $M$ est une classe mod $\underline{\mathrm{Pic}}^n_{X/k}$.

Le th. de \deleted{finitude} R.R. \inserted{(faible) sur $X$} montre que les polynômes de Hilbert des $X\in M$ sont les \unsure{ms\%}. Le th. de finitude de Chow (repris ds les schémas de Hilbert) montre donc que $M$ est quasi-compact.
% slanted
\margin{plus simplement, les degrés des diviseurs de $M$ restent fixés \unsure{---|\ldots}}

Donc

\uline{Corollaire} $\underline{\mathrm{Pic}}^n_{X/k}$ est qu. cpct \deleted{\unsure{un}}

Cela équivaut à

\uline{Corollaire} $\underline{\mathrm{Pic}}^n_{X/k}/\mathrm{Pic}^0_{X/k}$ est fini

ou encore

\uline{Corollaire} $\underline{\mathrm{Pic}}^n_{X/k}=\underline{\mathrm{Pic}}^\tau_{X/k}$, le groupe de torsion \deleted{\illegible{1}} de $\text{NérSév}_{X/k}$ est fini

La relation $\underline{\mathrm{Pic}}^{\unsure{n|-|\tau}}_{X/k}=\mathrm{Pic}^\tau_{X/k}$ signifie aussi

\uline{Corollaire} Soit $N=\underline{\mathrm{Pic}}_{X/k}/\underline{\mathrm{Pic}}^\tau_{X/k}$. Alors la forme d'intersection sur $N$ est non singulière.
% archive sheet number: 68

\scanpage{70}
Notons dans ce contexte le résultat plus profond de Igusa

\uline{Proposition}. $X$ comme ci-dessus, \deleted{\unsure{soit}} posons
\[
E(\overset{\inserted{\unsure{k|1}}}{X})=\sum_{i=0}^{4}(-1)^i b_i(X)
\]
avec
\[
\left\{\begin{aligned}
b_0(X)&=b_4(X)=1\\
b_1(X)&=b_3(X)=\dim\underline{\mathrm{Pic}}_X\\
E(X)&=c_2(X)=\text{self-intersection de la diagonale dans } X\times X.
\end{aligned}\right.
\]
ce qui détermine donc $b_2(X)$ de façon unique, \inserted{$\boxed{b_2=E-2+2p}$ où $p=\dim\underline{\mathrm{Pic}}_X$} Ceci dit, posons
\[
\rho=\operatorname{rang}(\underline{\text{NérSév}}_{X/k}),
\]
alors on a l'inégalité
\[
\boxed{\rho\leq b_2}
\]
% archive sheet number: 69

\scanpage{71}
\uline{Théorème de finitude sur Picard} (Commentaires

Kleimann prouve, pour $X/k$ \uline{projectif}, les critères \textcircled{\scriptsize viii} de $\tau$-équivalence à $0$, sans utiliser les fourbis délicats (i) (ii) \ldots. Pour récupérer le cas $X/k$ propre, il doit utiliser (i) (via lemme de Chow). Avec le même grain de sel, il prouve \textcircled{\scriptsize v} que $\underline{\mathrm{Pic}}^\tau_{X/S}$ est qu. cpct. sur $S$. [\uline{En fait, il prouve directement que $\underline{\mathrm{Pic}}^{\mathrm{num}}_{X/S}$ est quasi-compact\ldots.}].

Utilisant le th. de l'index pour une surface projective non singulière, Kleimann prouve de plus le ``critère d'équivalence'' sous la forme suivante : Si $X$ est projective $/k$, ttes \unsure{ses|les|es} composantes de dim $\geq3$, et $Y$ section hyperplane, d'où $f^*\colon\underline{\mathrm{Pic}}_{\deleted{\unsure{Y}}X/k}\unsure{-|\to}\underline{\mathrm{Pic}}_{Y/k}$, alors $f^{*^{-1}}(\underline{\mathrm{Pic}}^\tau_{Y/k})=\underline{\mathrm{Pic}}^\tau_{X/k}$, i.e. si $k$ \unsure{sp.|sep.|alg.} clos, $\mathrm{NSL}(X)\to\mathrm{NSL}(Y)$ est injectif. Ce\deleted{\unsure{ci}}\inserted{ qui précède}, et l'``inégalité'' de \uline{Picard. Igusa}, \inserted{et le th. de \uline{Néron-Sévéri} \deleted{\unsure{pour}} pour les surfaces proj. non sing., et \unsure{résolution} des surfaces \inserted{de \uline{Abhyankar}}} \unsure{suffisent|suffit} à entraîner \textcircled{\scriptsize vi} que si $X/S$ est propre de prés. finie sur \deleted{$k$} $S$ quasi-compact, les groupes de N.S. des fibres géométriques sont de \deleted{rang} \inserted{type fini} ``majoré'' \deleted{\unsure{et de}} (i.e. groupe de torsion d'ordre majoré, et rang majoré). [Dans le cas $X/S$ projectif, cela sort donc des dévissages de Kleimann, sans utiliser mes fourbis]. On récupère aussi ainsi \textcircled{\scriptsize vii}.

\uline{TSVP}
% archive sheet number: 70

\scanpage{72}
Ces résultats impliquent ils le ``critère \unsure{de'|d'} équivalence'' sous la forme (ii) ? Quitte à ``nettoyer'' la situation, on est ramené au cas où $\underline{\mathrm{Pic}}_{X/S}$, $\underline{\mathrm{Pic}}_{Y/S}$ sont représentables, $\underline{\mathrm{Pic}}^\tau_{X/S}$ et $\underline{\mathrm{Pic}}^\tau_{Y/S}$ aussi et plats et fermés dedans, et à examiner $\underline{\mathrm{NSL}}_{X/S}\to\underline{\mathrm{NSL}}_{Y/S}$. Or tout ce qu'on sait par ce qui précède, c'est que c'est un \uline{monomorphisme}. Il n'est pas du tout clair pourquoi il serait de type fini. En fait, ma démonstration prouve que ces morphismes $\underline{\mathrm{Pic}}_{X/S}\unsure{-|\to}\underline{\mathrm{Pic}}_{Y/S}$, $\underline{\mathrm{NSL}}_{X/S}\to\underline{\mathrm{NSL}}_{Y/S}$, sont génériquement au dessus de $S$ des morphismes quasi-affines (peut-être $\unsure{\overline{\mathrm{m}}|=}$ affines). Moyennant (i), on montre d'ailleurs que (ii) et (iii) sont essentiellement équivalents [le critère en termes du pol. de Hilbert, ou au choix, en terme de $D^2.H^{\unsure{n|N}-2}$ et $D.H^{n-1}$]. Il ne semble pas que cela sorte des \unsure{tapés|topés|topos} Kleimann.
% archive sheet number: 71

\scanpage{80}
\unsure{Extension|Extensions} de Serre d'une VA et faisceaux inversibles à connexion.
% archive sheet number: 79

\scanpage{81}
\diagram{\begin{tikzcd}[ampersand replacement=\&]
X\arrow[d,"f"]\\ S
\end{tikzcd}}
\deleted{\unsure{propre}} plat de prés. \deleted{\unsure{pro}} finie

\[
\underline{\mathrm{Div}}_{X/S}\simeq
\underline{R}^{*}_{X/S}/\underline{O}^{*}_{X\deleted{\unsure{/S}}}
\xrightarrow{\frac{df}{f}}\underline{R}_{X/S}(\underline{\Omega}^1_{X/S})/\underline{\Omega}^1_{X/S}
\tag*{(*)}
\]
\deleted{\unsure{extension}}

d'où une extension [des couples $(D,\omega)$ tels que \deleted{\unsure{$dD$}} $d\log D\equiv\omega\;(\underline{\Omega}^1_{X/S})$]
\[
0\to\underline{\Omega}^1_{X/S}\to\underline E_{X/S}\to\underline{\mathrm{Div}}_{X/S}\to0
\]
\margin{\uline{N.B}. L'hom (*) est injectif si $S$ de car. nulle et $f$ normal \deleted{les fibres géométriques \unsure{normales}} (\unsure{loc.|donc} de prés. finie\ldots.)}

et par suite une suite exacte
\diagram{\begin{tikzcd}[ampersand replacement=\&,column sep=small]
0\arrow[r]\&f_*(\underline{\Omega}^1_{X/S})\arrow[r]\arrow[d,equal]\&f_*(\underline E_{X/S})\arrow[r]\&f_*(\underline{\mathrm{Div}}_{X/S})\arrow[r,"\partial"]\&R^1 f_*(\underline{\Omega}^1_{X/S})\\
\&\underline{\omega}_{X/S}
\end{tikzcd}}
\margin{à vérifier que c'est l'homomorphisme bien connu déduit de $\underline{\mathrm{Pic}}_{X/S}\to R^1 f_*(\underline{\Omega}^1_{X/S})$. c'est trivial}

ou encore
\diagram{\begin{tikzcd}[ampersand replacement=\&,column sep=small]
0\arrow[r]\&\underline{\omega}_{X/S}\arrow[r]\&f_*(\underline E_{X/S})\arrow[r]\&f_*(\underline{\mathrm{Div}}_{X/S})^{\unsure{w|W}}\arrow[r]\arrow[d,equal]\&0\\
\&\&\&\operatorname{Ker}(\partial\colon f_*(\underline{\mathrm{Div}}_{X/S})\to R^1 f_*(\underline{\Omega}^1_{X/S}))
\end{tikzcd}}

Considérons
\[
\underline P_{X/S}=f_*(\underline R^*_{X/S})/f_*(\underline O^*_X)\hookrightarrow f_*(\underline{\mathrm{Div}}_{X/S})\unsure{\leftarrow|-}
\]
On a, du moins si $X\to S$ est \uuline{quasi-projectif} \deleted{\unsure{et}}
\[
f_*(\underline{\mathrm{Div}}_{X/S})/\underline P_{X/S}\simeq\underline{\mathrm{Pic}}_{X/S}
\]
\deleted{d'autre part, \unsure{on a}}

et on pose
\[
\begin{aligned}
f_*(\underline{\mathrm{Div}}_{X/S})^w/\underline P_{X/S}
&=\underline{\mathrm{Pic}}^w_{X/S}\subset\underline{\mathrm{Pic}}_{X/S}\\
&=\operatorname{Ker}(\underline{\mathrm{Pic}}_{X/S}\to R^1 f_*(\underline{\Omega}^1_{X/S}))
\end{aligned}
\]

$*$ On \unsure{adjoint|définit} la flèche oblique
% archive sheet number: 80

\scanpage{82}
\deleted{Massy le 24.4.1967}\\
\deleted{Monsieur P. Guetta}\\
\deleted{Mouvement du Milliard}\\

\deleted{Monsieur,}\\
% archive sheet number: 81

\scanpage{83}
\diagram{\begin{tikzcd}[ampersand replacement=\&,column sep=small]
\&\&\&\underline P_{X/S}\arrow[dl,"{\unsure{\hat d|\dot d|j}}",swap]\arrow[d]\\
0\arrow[r,dash]\&\underline{\omega}_{X/S}\arrow[r]\&f_*(\underline E_{X/S})\arrow[r]\&f_*(\underline{\mathrm{Div}}_{X/S})^{w}\arrow[r,dash]\&0
\end{tikzcd}}
définie par
\[
f_*(\underline R^{\unsure{\times|*}}_{X/S})\to f_*(E_{\deleted{\unsure{Y}}X/S})
\]
donnée par
\[
f\mapsto\deleted{\unsure{f}}\left(\mathrm{div}\,f,\frac{df}{f}\right)
\]

Il faut vérifier que c'est nul sur $f_*(\underline O^*_X)$, cela marche si \inserted{$f$ \uuline{propre et plat}, \deleted{$\underline O_X$} $\underline O_X$ \uuline{est \unsure{coh.} plat en dim $0$ sur $S$ et}} $f_*(\underline O_X)$ \uuline{est étale sur $S$} (p.ex. si $\mathcal O_S\simeq f_*(\underline O_X)$ universellement), ce que nous allons supposer.
\margin{En l'absence de cette hypothèse, on obtient seulement une ext. de $\underline{\mathrm{Pic}}^w_{X/S}$ par $\underline{\omega}_{X/S}/\operatorname{Im} f_*(\underline O_X^*)$}
\diagram{\begin{tikzcd}[ampersand replacement=\&,row sep=tiny,column sep=tiny]
\&X\arrow[dl]\\ S'\arrow[dr]\\ \&S
\end{tikzcd}}
Nous obtenons alors une extension
\[
\mathcal M_{X/S}=f_*(\underline E_{\deleted{\unsure{Y}}X/S})/\underline P_{X/S}
\]
\diagram{\begin{tikzcd}[ampersand replacement=\&,column sep=small]
0\arrow[r]\&\underline{\omega}_{X/S}\arrow[r]\arrow[d,dash]\&\mathcal M_{X/S}\arrow[r]\&\underline{\mathrm{Pic}}^w_{X/S}\arrow[r]\&0\\
\&\substack{\text{fibré vectoriel}\\ \text{sur }S}
\end{tikzcd}}

\uline{Cas favorable} $\underline{\Omega}^1_{X/S}$ \uuline{est \unsure{ch.|coh.} plat sur} $S$ \uuline{en} dim $0$, i.e. $\underline{\omega}_{X/S}$ est un fibré vectoriel associé à un module \uline{loc. libre} sur $S$. [C'est le cas si $f$ est lisse et projectif et $S$ de car nulle (?), ou si $X/S$ est un schéma abélien]. Alors $R^1 f_*(\underline{\Omega}^1_{X/S})$ est aussi représentable par un fibré vectoriel, donc si $\underline{\mathrm{Pic}}_{X/S}$ est représentable (p.ex. $f$ à fibres \unsure{géom.|géom} intègres) $\underline{\mathrm{Pic}}^w_{X/S}$ l'est aussi. Donc par descente $\mathcal M_{X/S}$ l'est également.
% archive sheet number: 82

\scanpage{84}
\deleted{Cher Monsieur Guetta,}\\
% archive sheet number: 83

\scanpage{85}
\noindent Définition \uline{directe} de l'extension $\underline{\mathcal M}_{X/S}$.

Soit $\mathcal M^\sharp_{X/S}$ le foncteur $(\mathrm{Sch})_{/S}^{\circ}\longrightarrow(\mathrm{Ens})$, faisceau \inserted{fppf} associé au préfaisceau
\[
\begin{aligned}
 &\deleted{\mathcal M'_{X/S}(T)=}\\
 Q:\quad T&\rightsquigarrow \text{classes, à isom près, de faisceaux}\\
 &\phantom{\rightsquigarrow{}}\text{inversibles sur }X\times_S T,\text{ munis d'une}\\
 &\phantom{\rightsquigarrow{}}\text{\uline{connexion} relativement à }T.
\end{aligned}
\]

\deleted{On a une} \deleted{N.B} On a une suite exacte de pré\uline{faisceaux} en $T$
\[
 H^0(X,\underline O_X^*)\xrightarrow{\varphi\mapsto\frac{d\varphi}{\varphi}}
 \frac{H^0(X,\Omega^1_X)}{\deleted{\unsure{\operatorname{Im}|d\log}}}
 \longrightarrow Q(T)\longrightarrow\mathrm{Pic}(X_T)\longrightarrow H^1(X,\underline\Omega^1_{X/S})
\]
d'où en passant aux faisceaux en $T$ associés
\[
 0\longrightarrow f_*(\underline\Omega^1_X)/\underbrace{\operatorname{Im}f_*(\underline O_X^*)}_{\text{par }\varphi\mapsto\frac{d\varphi}{\varphi}}
 \longrightarrow\mathcal M^\sharp_{X/S}\longrightarrow\underline{\mathrm{Pic}}^{w}_{X/S}\longrightarrow0
\]
où on a défini encore
\[
 \underline{\mathrm{Pic}}^w_{X/S}=\operatorname{Ker}\bigl(\underline{\mathrm{Pic}}_{X/S}\longrightarrow R^1 f_*(\underline\Omega^1_{X/S})\bigr)
\]

On va définir un homomorphisme d'ext.
\diagram{\begin{tikzcd}[ampersand replacement=\&,column sep=small]
0\arrow[r,dash]\&\deleted{\unsure{f_*|f}}\underline{\unsure{\omega|\Omega}}_{X/S}\arrow[r]\arrow[d,"{\mathrm{can}}"]\& f_*(\underline E_{X/S})\arrow[r,dash]\arrow[d,"\alpha"]\&\underline{\mathrm{Div}}^w_{X/S}\arrow[r,dash]\arrow[d,"{\mathrm{can}}"]\&0\\
0\arrow[r]\&\underline\omega_{X/S}/\operatorname{Im}f_*(\underline O_X^*)\arrow[r]\&\mathcal M^\sharp_{X/S}\arrow[r]\&\underline{\mathrm{Pic}}^w_{X/S}\arrow[r]\&0
\end{tikzcd}}

Pour ceci, on va interpréter \deleted{$\underline E_{X/S}$} l'extension $\underline E_{X/S}$ de $\underline{\mathrm{Div}}_{X/S}$ par $\underline\Omega^1_{X/S}$, en notant que pour un diviseur relatif donné $D$ sur un
% archive sheet number: 84

\scanpage{86}
\diagram{\resizebox{\linewidth}{!}{\begin{tikzpicture}[x=1cm,y=-1cm,every node/.style={inner sep=1pt}]
\node at (6.4,.7) {\begin{tikzcd}[ampersand replacement=\&,row sep=tiny,column sep=tiny] X'\arrow[d,dash]\\X\&\unsure{U|u|\mathcal U}\end{tikzcd}};
\node at (2.6,2.7) {$\begin{array}{ccc}P_i&P_{ij}&P_{ji}\\[4pt]X_i&\multicolumn{2}{c}{X_i\times_X X_j}\\[4pt]X&&\end{array}$};
\node at (6.3,3.7) {\begin{tikzcd}[ampersand replacement=\&,row sep=small]X'_\alpha\arrow[d,dash]\\ X\end{tikzcd}};
\node at (9.2,4.7) {\begin{tikzcd}[ampersand replacement=\&,row sep=small,column sep=small]G\&Z\arrow[d,"{\scriptscriptstyle\triangledown}"{description,pos=.12}]\\\& S\end{tikzcd}};
\node at (11.1,3.9) {$\deleted{\unsure{E|Z}}$};
\node at (12.6,4.8) {$Z\times^G E\simeq X$};
\node at (12.6,5.6) {$Z\times^G E'\simeq X$};
\node at (2.6,5.4) {$\deleted{\unsure{X_\alpha|G_\alpha}}$};
\node at (4.8,6.1) {\begin{tikzcd}[ampersand replacement=\&,column sep=small,row sep=small] Z_\alpha\arrow[d]\&X'\arrow[dl]\\X\&{}\end{tikzcd}};
\node at (8.55,6.4) {$\deleted{\unsure{g|G}}$};
\node at (12.1,6.7) {$E\longrightarrow E'$};
\node at (9.5,7.7) {$Z\times^G E$};
\node at (12,7.7) {$G$};
\node at (2.4,8.1) {$Z_\alpha$};
\node at (4.3,8.1) {$\deleted{\unsure{g|G}}$};
\node at (2.8,9.3) {$Z_\alpha\times^{G_\alpha}_S H$};
\node at (6.3,9.3) {$Z_\alpha\times^{G_\alpha}_S E$};
\node at (8.1,11.05) {$\operatorname{Hom}_{\unsure{G_\alpha|G\times\alpha}}(E,F)\xrightarrow{\sim}\operatorname{Hom}_{S,\deleted{\unsure{G_\alpha|Z_\alpha}}}(T(E),T(F))$};
\node at (13.5,11.5) {$\unsure{380|3.80}$};
\node at (2.1,11.9) {$S$};
\node at (5.7,12.4) {$(\mathrm i)\quad \deleted{\unsure{E|Z}}Z_\alpha$ connexe};
\node at (9.5,13.6) {\begin{tikzcd}[ampersand replacement=\&,column sep=0pt,row sep=small]
{(\mathrm{ii})\quad\operatorname{Hom}(}\& {Z_\alpha/H,}\&{Z/K)}\&{\xleftarrow{\sim}\operatorname{Hom}_{\unsure{C|G}}(G/H,G/K)}\\
{}\&Z\arrow[u]\arrow[ur]\&{}\&{}\end{tikzcd}};
\node at (13.1,14.3) {$g\in G/K,\quad gKg^{-1}\supset H$};
\node at (6,15.3) {$Z\longrightarrow Z/K$};
\node at (6.4,16.2) {$\operatorname{Hom}_S(Z,Z/K)\simeq G/K$};
\node at (11.6,15.8) {$Z\mathbin{\times}^{G}H$};
\node[align=left] at (14.1,16.4) {\deleted{\unsure{les|un}}\\\unsure{tout|tous|est} \unsure{autom|automs|autour}\\de $Z\times^G H$\\\unsure{provient|provient de|principal} loc.\\d'un};
\node at (1.25,17.65) {\begin{tikzcd}[ampersand replacement=\&,row sep=small]
{}\arrow[d,dashed,dash]\\ Z_{2}^{\unsure{4|\sharp}}\arrow[d]\\ Z_1\arrow[d,leftrightarrow]\\ Z
\end{tikzcd}};
\node at (10.7,17.6) {$Z$ connexe};
\node[draw,rounded corners=2pt] at (12.2,18.25) {$Z/H$\quad tout \unsure{autom|auton|autour} descend de $G$.};
\draw[thin] (2.28,1.98) to[bend right=3] (10.62,18.75);
\draw[thin] (10.34,.02) to[bend right=2] (13.1,17.1);
\draw[thin,->] (.05,11.4) to[bend right=5] (6.3,19.65);
\end{tikzpicture}}}
% archive sheet number: 85

